
\begin{quote}
Every quantitative claim below is established in Chapters 4--6 and reproducible from the repository; this chapter interprets, connects and stress-tests them.
\end{quote}

This chapter discusses the results rather than reporting new ones: what the project achieved against its two research questions, what the remaining failures are made of, what the measurements say about the dataset's own annotation process, and, in \S7.7, an answer to each of the five objections \S2.9 raised against the approach before any result was known. The objective-by-objective audit belongs with the conclusions and is in \S9.1.

\section{Achievement against the research questions}

\textbf{RQ1} asked whether spatial-relationship annotation can be automated at a quality comparable to human annotation. The answer is predicate-shaped rather than singular. For the lateral and proximity predicates the tool is, by every measure available, \textit{at least} as good as the human process: 0.97--0.998 recall, audited true precision $\approx$ 1.0, and for \texttt{near} a perfect held-out score against the one annotator who used the label and never influenced the threshold. For support, after two evidence upgrades motivated by measurement (depth co-location, mask contact), recall is 0.88/0.81 with audited precision around 0.9, comfortably comparable. For the depth pair the cascade of relative depth and the ground-plane fallback reaches 0.64/0.66 pooled (0.84 once the two inverted-convention groups are aligned), and where the tool commits it agrees at 0.95--1.00 with six of the seven same-convention annotators, the seventh being the dataset's smallest sample at 65 triplets; the remaining shortfall is calibrated abstention plus that inverted direction convention. "Comparable to human quality" understates that case: on front/behind the tool is more consistent than the human process it is measured against.

\textbf{RQ2} asked whether the automatic labels can train a relation model as well as human labels. The controlled experiment answered more strongly than the question was posed: with identical features, model, seed and split, the auto-trained classifier reaches 0.76 mean recall against held-out \textit{human} gold versus 0.30 for its human-trained twin. At this dataset's annotation scale, the automatic labels are better training material than the labels they were validated against. The mechanism is not mysterious: density (20$\times$ more triplets) and consistency (one definition, uniformly applied). But it converts the project's premise from "removing the bottleneck loses little" to "removing the bottleneck gains".

The third arm is what makes that claim hard to dismiss. Self-training on the human labels, the standard semi-supervised remedy for exactly this problem, reaches 0.36 and closes only 15\% of the gap, and its bookkeeping shows why: the teacher contributes about a thousand confident \textit{negative} pseudo-labels for every positive one, propagating the annotators' silence rather than their judgement, and on \texttt{near} it drives recall below the human baseline it started from. The comparison therefore is not against doing nothing, but against the obvious alternative, under identical conditions.

The six objectives of \S1.2.1 are audited against their evidence in \S9.1; this chapter is concerned with what the results \textit{mean} rather than with whether each objective was discharged.

\section{What the remaining failures are made of}

The failure gallery diagnoses every one of the 1,689 missed human triplets by re-checking rule conditions, so the failure analysis is exhaustive rather than anecdotal. Three observations matter most.

First, \textbf{genuine tool error is rare}: depth-ordering mistakes are 1--5\% of front/behind misses; across all predicates, misses attributable to avoidable error are roughly 7\%. The bulk of the miss mass is calibrated abstention (depth ambiguity band, 52--61\% of front/behind misses) and measured annotator defects (38--42\%, a share that \textit{grew} as the ground-plane fallback shrank the abstention share around it).

Second, \textbf{the support arc shows the method working as a method}. The box rule shipped at \textasciitilde{}0.27 true precision; the audit localised the failure (projection adjacency), the gallery localised the misses (containment), one geometric insight (stacked objects share a camera distance) fixed half the false fires, and one perception upgrade (mask-bottom contact) fixed most of the rest while \textit{raising} recall, each step calibrated on train annotators and validated held-out. The residual failure mode is precisely characterised too: a person \textit{holding} an object satisfies pixel contact (3 of the 7 remaining audited errors), and a class-aware guard is the documented next refinement.

Third, \textbf{what looked like a depth-resolution ceiling was mostly a rules ceiling, and it moved}. Two objects at similar camera distance cannot be ordered by relative monocular depth (the \texttt{depth\_eps} sweep bounds the trade: recall up to \textasciitilde{}0.71 at $\epsilon$=0 for \textasciitilde{}0.26--0.36 precision), but the ground-plane fallback recovers most of the abstention band \textit{without} depth, from pure projection, once the tool's own contact evidence guards against elevated objects. What is left is narrower and equally well characterised: objects resting on supports the detector has no box for, and pairs whose bottom edges tie within the band. Both operating points are documented, revisable decisions rather than hidden constants (ablations A2, A7).

Section 4.12 settles what the ablations could not. If the front/behind gap were depth \textit{noise}, estimates jittering either side of a boundary, moving the camera would flip verdicts, since that is the perturbation such noise responds to. It does not: the predicate reproduces itself 0.955 of the time, above \texttt{on} and \texttt{under}, and holds at 0.924 at 89-fold compression. A predicate recovering 0.64 of the human labels while agreeing with itself at that rate is not guessing; it is applying a consistent criterion the annotators did not share. With A8, where quadrupling the depth model changed nothing, and \S4.5, where two groups labelled the pair oppositely, the weight of the shortfall sits on definitional disagreement rather than perception. That does not dissolve monocular ambiguity, which genuinely bounds the predicate at equal camera distance, but it relocates most of the measured gap away from it. The implication is unglamorous: the best expected return here is a written annotation guideline, not a better depth network.

\section{The dataset's annotation process, examined}

The source paper flagged \texttt{near} as inconsistent and called for "clear annotation guidelines (e.g., spatial thresholds for 'near')". This project quantifies how much further the guideline problem goes:

\begin{enumerate}
\item \texttt{near} was used by 3 of 9 annotator groups, with \textasciitilde{}4$\times$ variance in how
exhaustively equally-close pairs were labelled. Yet all three annotators' labels sit inside one fitted threshold (held-out recall 1.0): consistent \textit{notion}, non-exhaustive \textit{application}.
\item Two annotator groups applied the \textbf{inverted direction convention} for
in front of / behind (2--5\% agreement where the tool commits; flipping recovers 0.93/0.71).
\item Support pairs were often labelled in \textbf{one direction only} (one group
all-\texttt{on}, another all-\texttt{under}).
\item The official guidance, confirmed at the annotation tool's repository, is
\textbf{vocabulary lists with no definitions} (\S4.7), which is what makes every defect above predictable rather than surprising.
\end{enumerate}

This reframes the evaluation itself: for several predicates there is no human consensus to agree with, only per-annotator behaviours. The dissertation's response, per-annotator reporting, annotator-aware calibration and operational definitions as the deliverable, is to our knowledge the first time this dataset's label semantics have been made explicit.

The "tenth annotator" framing survives contact with the data, and \S4.6 puts numbers on it. The tool is deterministic, the same labeller for every group, so the 0.216 spread in its agreement across the seven consistent annotators (0.717 to 0.933) measures their heterogeneity, not its inconsistency. Fréchet bounds with the tool as common reference place annotator-to-annotator agreement in [0.74, 0.92], containing the tool's own 0.869. The claim is deliberately modest, since the bounds assume the batches are exchangeable: the automatic annotator cannot be shown to agree with the humans any less well than they can be shown to agree with each other. Without overlapping assignments the quantity cannot be measured outright, and that absence is itself a finding about the dataset's construction, one a replication should design away by having two annotators share a batch.

\section{Methodological reflection}

Choices that proved right: the \textbf{PredCls isolation} (without it, every rule result would be confounded by detection; the SGDet decomposition shows the relation layer at 0.85 conditional mean, invisible in the 0.38 end-to-end number); the \textbf{sparse-gold protocol} (recall-primary plus restricted precision plus audits; the audit overturned the naive reading of restricted precision for five predicates and confirmed it for support); and \textbf{train-only calibration with held-out annotators} (every fitted threshold generalised: near recall 1.0, support F1 0.87 on annotators the thresholds never saw).

One choice proved its worth only afterwards. The withdrawn single-seed claim of \S6.3.1 shows the discipline applied faithfully to thresholds, fitting on some annotators and validating on others, reaching \textit{model training variance} late. The fix was cheap, four extra runs on a free GPU tier, and a replication designed from the start would have trained every arm at three seeds and reported ranges throughout, which is what the final version does.

Choices a stricter replication should improve: the \textbf{audits were verdicted by the author} (conservatively, with verdicts and rendered evidence published for spot-checking; the independent validation study of Appendix E.3 is the designed remedy, and blind external verdicting should have been the instrument from the first audit rather than the last); the \textbf{support-rule iteration used the same audit machinery twice}, so the second audit is confirmatory rather than fully independent; the SGDet \textbf{prompt/ threshold tuning used one disclosed iteration on a trial slice} that over-estimated full-set detection quality, a small, instructive example of trial-set optimism; and 2,000-scene invariant fuzzing pins rule consistency but not rule \textit{truth}, which only the audits address.

\section{Synthesis against the geometry-to-label lineage}

The pipeline borrows its skeleton from the SpatialVLM family (Chen et al., 2024) and, further back, from CLEVR (Johnson et al., 2017); neither supplies the recovery of geometry from real photographs, where this project's difficulty sits. What it adds is not the skeleton but the parts the lineage leaves implicit, because SpatialVLM and VQASynth (Remyx AI, 2024) generate \textit{training text} at internet scale without ever confronting a fixed predicate vocabulary with human ground truth, and SpatialRGPT (Cheng et al., 2024) curates region representations with depth but validates downstream rather than against annotators. This project's contributions are \textbf{annotator-aware calibration} (fit only on annotators who used a label; hold out annotators, not just images), \textbf{contact as the support signature} (mask-bottom adjacency, unused by the box-geometry lineage, which repaired both error directions at once and parallels the argument for pixel-accurate grounding in panoptic scene-graph generation (Yang, J. et al., 2022)), and \textbf{loss attribution as methodology} (every miss diagnosed; every gap decomposed into abstention against annotator against error, detection against relations). RoboSpatial's reference-frame taxonomy (Song et al., 2025), cited in Chapter 3 to justify camera-frame laterality, proved the right lens for a \textit{measured} phenomenon: the front/behind inversion is a reference-frame disagreement inside one annotation team, an instance of the frame-dependence Landau and Jackendoff (1993) describe.

Two results connect this project to literatures outside its lineage. Chapter 6's benchmark finding is a case of the problem Northcutt, Athalye and Mueller (2021) demonstrated across ten benchmarks, that errors in test annotation change model rankings and can select the wrong model; here the defect is systematic annotator convention rather than random noise, and it distorts a ranking between two label \textit{sources} rather than two architectures. The diagnosis matches the motivation for SpatialSense (Yang, K., Russakovsky and Deng, 2019) and Rel3D (Goyal et al., 2020), both built after their authors found relation benchmarks could be scored well without using spatial information; this dissertation observes the mirror image, a model scoring well by reproducing annotator selection habits. The RQ2 result is Ratner et al.'s (2017) weak-supervision prediction confirmed in a domain the original work did not address, with the standard semi-supervised alternative (Lee, 2013) implemented and measured rather than argued away.

\section{Limitations and threats to validity}

\textbf{Internal.} Thresholds are fitted to six annotator groups of one dataset; audits are author-verdicted (\S7.4); the two-stage audit shares machinery with the rule change it evaluates.

\textbf{External.} One laboratory domain, six object classes, one camera and mounting; the fitted constants (\texttt{near\_T}, $\epsilon$ values, contact threshold) are dataset-specific by design, and the transferable artefact is the \textit{procedure}, fit on some annotators and validate on held-out ones, not the numbers. The only out-of-domain evidence is qualitative (Appendix E.4) and measures nothing, since no labelled out-of-domain gold exists; a modest labelled cross-domain sample is the missing experiment, cheap enough that a replication should simply include one. Full automation is currently detection-bounded (0.38 end-to-end with a worst-case zero-shot detector; the authors' trained detector would close most of that gap, unverified without their weights).

\textbf{Construct.} "Agreement with human labels" is an imperfect proxy when the humans disagree with each other; the per-annotator decompositions mitigate but cannot eliminate this. The RQ2 result compares supervision \textit{at this dataset's annotation scale}. It does not claim automatic labels beat abundant, guideline-driven human annotation, a regime this dataset does not contain.

\textbf{Ethics.} Scene images contain identifiable people; figures for publication use the dataset as released (CC-BY 4.0) with faces blurred as a courtesy.

\section{The objections of \S2.9, answered}

The literature review stated five objections to a rule-based annotator before any result was reported, so that this chapter could be read as an attempt on them. Each is answered here with the evidence that bears on it, including the two the evidence does not settle.

\textbf{Vocabulary scale: not answered, and conceded.} Nothing in this dissertation bears on predicates beyond the seven, and the objection is that a hand-written rule set does not extend by learning. It stands in full. What the work can offer is a boundary rather than a rebuttal: the rules are decidable because the predicates are spatial, and \S3.3 makes that dependence explicit rather than hoping it generalises. Section 9.3 records it as the limitation most likely to matter to anyone reusing the method.

\textbf{Systematic error: partly confirmed, partly refuted, and the split is informative.} The objection predicts that a downstream model absorbs the rule's blind spot as though it were fact. Chapter 6 is that prediction coming true and is reported as such: the auto-trained arm ranks below the human-trained one on the metric the field uses, and the gap is concentrated where the annotation itself is defective. But the objection also predicts that consistent-but-wrong supervision is worse than inconsistent human supervision, and on this dataset that is refuted at every level where the question was asked, by a factor of two and a half in the controlled experiment and by a planner that clears the occluder in 22 of 25 scenes against 0 with no relations at all. The reconciliation is that systematic error is worse than random error only when it is \textit{wrong}; \S4.12's finding that front/behind agrees with itself across viewpoints 0.955 of the time while agreeing with the annotators 0.64 of the time is the shape of a consistent rule meeting a different convention, not of a consistent mistake.

\textbf{Circular validation: conceded.} Section 7.4 already lists this first among the choices a stricter replication should improve, and no result in Chapters 4 to 6 removes it. What belongs here is the extent of the concession rather than a restatement of it. The circularity bounds the \textit{precision} estimates, because those rest on verdicts the author gave. It does not reach the structural guarantees of \S3.6, which are checkable without any verdict at all, nor the downstream findings of Chapters 5 and 6, which are scored against human annotation the author did not produce and in which the automatic arm is judged by its rival's yardstick. The designed remedy is the study of Appendix E.3 and it has no results at the time of writing, which leaves this the objection a reader should weight most heavily.

\textbf{Reference frame: answered as far as it can be, which is not all the way.} The dissertation does not assert that the camera frame is correct; it shows that the tool applies one frame consistently (\S4.12), that two annotator groups applied another (\S4.5), and that recall rises from 0.64/0.66 to 0.84 once the convention is aligned. That establishes disagreement rather than error, and identifies which party is consistent. It does not establish which convention a robot should obey, and no measurement in this dissertation could, because that is a question about the specification and not about the data.

\textbf{Better collection rather than cheaper labels: not attempted, and partly answered sideways.} The project did not re-collect anything. It did, however, produce the written operational definitions the annotation process never had (\S3.5) and quantify three specific defects those definitions would have prevented, which is the input a better collection round would need. Cheap labels and correct definitions are complements rather than alternatives, and this work supplies evidence for the second while delivering the first.

\section{Aims, revisited}

The project set out to show that automatic labels are \textit{not much worse} than human ones and found conditions under which they are decisively better for downstream learning, those conditions being sparse, guideline-free annotation: exactly the regime the source dataset occupies. That reframes what the bottleneck was doing. It was not only slowing the dataset down, it was limiting what the dataset could teach.